%%═══════════════════════════════════════════════════════════════════
%% Lecture 12 — Advanced ODE Methods
%% 77315 — Computational Physics
%%═══════════════════════════════════════════════════════════════════

\section{Lecture 12 — Advanced ODE Methods}\label{sec:lec12}
\hrule\vspace{1.5em}

%%──────────────────────────────────────────────────────────────────
\subsection*{Overview}
Building on Euler's method, we develop higher-order integrators for
initial-value problems $\dot{y}=f(x,y)$.  The main families are:
\emph{Adams multistep} methods (Adams–Bashforth explicit,
Adams–Moulton implicit, predictor–corrector combinations) and
\emph{Runge–Kutta} single-step methods (RK2, RK3, RK4).
We also discuss adaptive step-size control via Richardson
extrapolation and the treatment of stiff ODEs.

%%──────────────────────────────────────────────────────────────────
\subsection{Adams--Bashforth (Explicit Multistep)}\label{subsec:lec12:ab}

Adams–Bashforth methods use past values of $f$ to extrapolate.
The 2-step and 4-step formulae are

\begin{align}
  y_{n+1} &= y_n + h\left[\tfrac{3}{2}f_n - \tfrac{1}{2}f_{n-1}\right]
  \quad \bigO{h^3},
  \label{eq:lec12:ab2}\\
  y_{n+1} &= y_n + \frac{h}{24}\left[55f_n - 59f_{n-1}
              + 37f_{n-2} - 9f_{n-3}\right]
  \quad \bigO{h^5}.
  \label{eq:lec12:ab4}
\end{align}

The 4-step method requires only \emph{one} new function evaluation
per step (the previous three values are reused).  Starting values
$y_1,y_2,y_3$ must be obtained by a Runge–Kutta method.

%%──────────────────────────────────────────────────────────────────
\subsection{Adams--Moulton (Implicit Multistep)}\label{subsec:lec12:am}

The implicit Adams–Moulton formula uses $f_{n+1}$ (the unknown):
\begin{equation}
  y_{n+1} = y_n + \frac{h}{24}\left[9f_{n+1} + 19f_n
             - 5f_{n-1} + f_{n-2}\right]
  \quad \bigO{h^5}.
  \label{eq:lec12:am4}
\end{equation}
Because $f_{n+1}$ depends on $y_{n+1}$, \eqref{eq:lec12:am4} is an
implicit equation.  For non-stiff problems it is solved by the
\emph{predictor–corrector} approach.

%%──────────────────────────────────────────────────────────────────
\subsection{Predictor--Corrector}\label{subsec:lec12:pc}

\begin{enumerate}
  \item \textbf{Predict} $y_{n+1}^{(0)}$ with Adams–Bashforth~\eqref{eq:lec12:ab4}.
  \item Evaluate $f_{n+1}^{(0)} = f(x_{n+1}, y_{n+1}^{(0)})$.
  \item \textbf{Correct} with Adams–Moulton~\eqref{eq:lec12:am4}.
  \item (Optionally iterate correction once more.)
\end{enumerate}

The combination achieves 5th-order accuracy with four function
evaluations per step (same as the four-step AB formula alone), but
the implicit corrector dramatically improves stability.

%%──────────────────────────────────────────────────────────────────
\subsection{Runge--Kutta Methods}\label{subsec:lec12:rk}

RK methods are single-step: no past values needed.

\subsubsection*{RK2 (Midpoint method)}
\begin{equation}
  k_1 = hf(x_n, y_n),\quad
  k_2 = hf\!\left(x_n+\tfrac{h}{2},\,y_n+\tfrac{k_1}{2}\right),\quad
  y_{n+1} = y_n + k_2 + \bigO{h^3}.
  \label{eq:lec12:rk2}
\end{equation}

\subsubsection*{RK4 (Classical)}
\begin{align}
  k_1 &= hf(x_n,\,y_n),\notag\\
  k_2 &= hf(x_n+\tfrac{h}{2},\,y_n+\tfrac{k_1}{2}),\notag\\
  k_3 &= hf(x_n+\tfrac{h}{2},\,y_n+\tfrac{k_2}{2}),\notag\\
  k_4 &= hf(x_n+h,\,y_n+k_3),\notag\\
  y_{n+1} &= y_n + \tfrac{1}{6}(k_1+2k_2+2k_3+k_4) + \bigO{h^5}.
  \label{eq:lec12:rk4}
\end{align}

RK4 costs 4 function evaluations per step and achieves
$\bigO{h^5}$ local truncation error.

%%──────────────────────────────────────────────────────────────────
\subsection{Adaptive Step Size and Richardson Extrapolation}
\label{subsec:lec12:adaptive}

\subsubsection*{Step-doubling / Richardson extrapolation}
Take one step of size $h$ to get $y^{(h)}$, then two steps of size
$h/2$ to get $y^{(h/2)}$.  For a method of order $p$:
\begin{equation}
  \text{error} \approx \frac{y^{(h/2)}-y^{(h)}}{2^p - 1}.
  \label{eq:lec12:rich}
\end{equation}
If the error exceeds the tolerance, halve $h$; if it is much smaller
than the tolerance, double $h$.

%%──────────────────────────────────────────────────────────────────
\subsection{Stiff ODEs}\label{subsec:lec12:stiff}

A system is \emph{stiff} when its Jacobian has eigenvalues with widely
separated magnitudes.  Explicit methods (RK4, Adams–Bashforth) require
$h \lesssim 1/|\lambda_{\max}|$, which may be extremely small even though
the solution of interest evolves on the slow timescale $1/|\lambda_{\min}|$.

\subsubsection*{Implicit and semi-implicit schemes}
Fully implicit methods (e.g.\ backward Euler, fully implicit
Adams–Moulton) have much larger stability regions and can handle
stiff problems with step sizes controlled by accuracy, not stability.
For systems where $f$ can be split as $f = f_{\rm fast} + f_{\rm slow}$,
\emph{semi-implicit} (IMEX) schemes treat only the fast part implicitly.

\nt{Stiffness is common in chemistry (fast reaction rates), nuclear physics, and coupled oscillators.  Always check whether an explicit integrator requires an unreasonably small $h$ before declaring it ``inaccurate''.}
